% appendices/prompts.tex
% LLM prompts used in the qualitative analysis of the expert survey free-text responses.

\chapter{LLM Prompts}
\label{ch:prompts}

All analyses used \texttt{gpt-4o-mini} via the OpenAI API with  \texttt{temperature=0.1} (structured JSON) or \texttt{0.2} (free-text).
Placeholders in braces are filled programmatically at runtime.

\paragraph{Data transmitted to the API and privacy safeguards.}
Only participants' open-ended (free-text) responses were sent to the
GPT-4o-mini endpoint, each accompanied solely by an internal reference
number (the CSV row index, shown as \texttt{[R\{rowId\}]}) so that
model-generated quotes could be traced back to the source data. No names,
affiliations, profile links, email addresses, demographic fields, or
closed-item responses were included in the API payload. All direct and
quasi-identifying fields---name, affiliation, and profile link---were
removed prior to any third-party processing, consistent with the
anonymization procedure documented in the codebook; the free-text
answers were additionally checked for embedded identifiers before
transmission. Only responses from participants who provided explicit
consent were included: non-consenting responses were removed by a consent
filter before any processing, including the LLM-assisted step. The consent
instrument in force throughout data collection explicitly informed
participants that ``open-ended responses may be analyzed using third-party
AI services (e.g., OpenAI); only de-identified text is transmitted and it
is not used to train the providers' models,'' on the legal basis of
explicit consent (GDPR Art.~6(1)(a)). Transmitted text was not used to
train the provider's models.

% ── Prompt set 1: per-question structured disagreement analysis ──────────────
\subsection{Per-Question Disagreement Analysis}
\label{app:prompt-per-question}

Used to produce the structured disagreement objects (one call per survey question per grouping dimension).

\paragraph{System}
\begin{quote}\small\ttfamily
You are a qualitative research analyst for an academic paper.
Your output will be used directly by the authors to identify specific
evidence for claims---not to write prose for them.
Be precise and quote actual phrases.
When citing representative\_quotes, include the [R\{rowId\}] citation
prefix exactly as it appears in the input.
Return ONLY valid JSON.
\end{quote}

\paragraph{User}
\begin{quote}\small\ttfamily
Survey question: ``\{question\_label\}''

Responses grouped by professional background
(each prefixed with [R\{rowId\}] = source CSV row):

\{grouped\_responses\}

Identify agreements and disagreements across groups.
Return ONLY valid JSON matching this schema exactly:

\{
  "agreements": ["<shared view across groups>"],
  "disagreements": [
    \{
      "topic": "<short label>",
      "positions": \{"<GroupName>": "<their view, or null if not expressed>"\},
      "representative\_quotes": \{
        "<GroupName>": "<[R\{rowId\}] short verbatim or near-verbatim phrase>"
      \}
    \}
  ],
  "group\_unique\_concerns": \{
    "<GroupName>": ["<concern mentioned only by this group>"]
  \}
\}
\end{quote}

% ── Prompt set 2: overall cross-cutting synthesis ────────────────────────────
\subsection{Overall Cross-Cutting Synthesis}
\label{app:prompt-overall}

Used once per grouping dimension after all per-question calls, to
produce the \texttt{overall} section of each JSON file.

\paragraph{System}
\begin{quote}\small\ttfamily
You are a qualitative research analyst for an academic paper on
AI-generated disinformation. Provide structured findings that the
authors can reference and cite.
When citing examples, include the [R\{rowId\}] citation prefix exactly
as it appears in the input.
Return ONLY valid JSON.
\end{quote}

\paragraph{User}
\begin{quote}\small\ttfamily
Below are selected open-ended survey responses from different groups in
an expert survey on AI-generated disinformation threats.
Each response is prefixed with [R\{rowId\}] = its source CSV row.

\{grouped\_responses\}

Identify cross-cutting agreements and disagreements.
Return ONLY valid JSON matching this schema exactly:

\{
  "cross\_cutting\_agreements": ["<shared view recurring across questions>"],
  "cross\_cutting\_disagreements": [
    \{
      "topic": "<short label>",
      "positions": \{"<GroupName>": "<their view>"\},
      "significance": "<why this disagreement matters for the paper>"
    \}
  ],
  "group\_profiles": \{
    "<GroupName>": "<2--3 sentence profile of their overall perspective>"
  \},
  "key\_tensions": [
    \{
      "tension": "<description>",
      "groups\_involved": ["<GroupName>"],
      "example": "<[R\{rowId\}] brief illustrative example from the data>"
    \}
  ]
\}
\end{quote}

% ── Prompt set 3: per-question free-text theme extraction ────────────────────
\subsection{Per-Question Theme Extraction}
\label{app:prompt-themes}

Used in \texttt{cluster\_open\_responses.py} to produce a plain-text
thematic summary per question (all respondents pooled).

\paragraph{System}
\begin{quote}\small\ttfamily
You are a qualitative research analyst.
Identify the main themes and key points from survey responses concisely.
Use bullet points. Be specific; quote short phrases and cite the
[R\{rowId\}] source identifier when referencing a specific response.
\end{quote}

\paragraph{User}
\begin{quote}\small\ttfamily
Survey question: ``\{question\_label\}''

Responses (n=\{N\}, each prefixed [R\{rowId\}] = CSV row number):
\{responses\}

List the 3--5 main themes/points raised across these responses.
\end{quote}

% ── Prompt set 4: cross-question overall synthesis ───────────────────────────
\subsection{Cross-Question Synthesis}
\label{app:prompt-synthesis}

\paragraph{System}
\begin{quote}\small\ttfamily
You are a qualitative research analyst summarising expert survey data
about AI-generated disinformation. Be concise and insightful.
\end{quote}

\paragraph{User}
\begin{quote}\small\ttfamily
Below are selected responses across all open-ended questions from an
expert survey on AI-generated disinformation threats.

\{responses\_by\_question\}

Provide a cross-cutting synthesis: What are the 5--7 overarching themes
or concerns that recur across multiple questions?
Note any tensions or disagreements.
\end{quote}

% ── Prompt set 5: profession group comparison ────────────────────────────────
\subsection{Profession-Group Comparison}
\label{app:prompt-profession}

\paragraph{System}
\begin{quote}\small\ttfamily
You are a qualitative research analyst.
Compare perspectives across professional groups concisely.
\end{quote}

\paragraph{User}
\begin{quote}\small\ttfamily
Below are open-ended survey responses grouped by the respondent's
professional background, from an expert survey on AI-generated
disinformation.

\{grouped\_responses\}

Identify:
1. Points where the groups broadly AGREE.
2. Points where the groups DISAGREE or have notably different emphases.
3. Any unique concerns that appear only within a specific group.
\end{quote}